In this section, we make the following assumptions about the strategy spaces and the utility functions of a network game $\calG$ defined by \eqref{def:game}.
\begin{enumerate}
    \item Every compact strategy space $X_i\subseteq \R^{m_i}$ is a \emph{basic semialgebraic set} $\calS(\bg_i)$ where each $g_{ik}$ is a real polynomial in $m_i$ variables and $K_i\neq\emptyset$ is a finite index set. We also assume the standard archimedean certificate of compactness  (\cref{def:archimedean}) on the sets.
    \item In every two-player game involving players $i,j\in N$, the utility functions $u_{ij}$ and $u_{ji}$ of player $i$ and $j$, respectively, are real polynomials in $m_i+m_j$ indeterminates.  Therefore, each utility function $u_i$ given by (\ref{def:pi}) is a real polynomial in $m_1+\dots + m_n$ indeterminates.
    \item The game $\calG$ is zero-sum.
\end{enumerate}
These three assumptions define a \emph{polynomial zero-sum network game} $\calG$. Glicksberg's theorem~(\ref{thm:Glick}) can be significantly strengthened for polynomial games --- each polynomial game has NE in which every player mixes only among finitely many pure strategies \cite{DresherKarlinShapley50,stein.ozdaglar.ea:separable}.

The solutions for polynomial zero-sum network games can be solved using Lasserre's hierarchies as described in \cref{sec:momsos}. Specifically, we relax the constraints on measures and polynomial inequalities of Program~\ref{P}. Instead of requiring that the mixed strategies of players lie in the cone of probability measures, we apply constraints on truncated sequences of moments. Similarly, we replace the inequalities \eqref{P:ineq} with the constraint that $w_i - U_i(\cdot, \bmu_{-i})$ lies in the truncated quadratic module.

The $h$-th level of the hierarchy is the following semidefinite program:
\begin{equation}
    \begin{aligned}
        \minimize_{\hat{\by}_1,\dots,\hat{\by}_n,\bw}\quad& w_1+\dots +w_n \\
        \text{\normalfont subject to}\quad
        & w_i - \bar{u}_i(\hat{\by}_{-i}) \in Q_{2h}(X_i) &\forall i \in N\\
        & \hat{\by}_i=(y_{i,\balpha})_{\balpha\in\N^{m_i}_{2h}} \text{ and $y_{i,\mathbf{0}}=1$} &\forall i \in N\\
        & M_h(\hat{\by}_i)\succeq 0 &\forall i \in N\\
        & M_h(g_{ik}, \hat{\by}_i)\succeq 0 & \forall k\in K_i,\, \forall i \in N\\
        & \bw=(w_1,\dots,w_n) \in\R^n \\
    \end{aligned}
    \label{momsdp}
\end{equation}

We then iteratively solve each level of the hierarchy until all sequences $\hat{\by}_i$ have a representing atomic measure, which we extract using the algorithm of \textcite[Algorithm 6.9]{Lasserre15:book}. With the strategies of other players fixed, the best response condition is equivalent to certifying optimality for a polynomial over a semialgebraic set.